% SOURCE_AUTHORITY: sga5_fr_workpass.tex, lines 14384--14403 (LF-normalized unit).
% SOURCE_SHA256: 791F4EFFC5E02832D5D77ED03518C8156D6F07E4C8238B03545DB93D883FBB28
Demonstremos a). O composto $\pi_{X/S}\circ\Fr_{X/S}=\fr_X$ é integral, sobrejetivo e radicial de maneira evidente; resulta que $\Fr_{X/S}$ é radicial (EGA I 3.5.6 (ii)); além disso, $\pi_{X/S}$ é separado e radicial, pois deduz-se de $\fr_S$ mediante a mudança de base $X\to S$; resulta que $\Fr_{X/S}$ é integral (EGA II 6.1.5 (v)) e sobrejetivo.

A demonstração de b) é evidente, observando que
\[
\JJ\OO_{(S,\fr_S)}=\JJ^p
\]
para todo ideal $\JJ$ de $\OO_S$.

Demonstração de c). Sabe-se que, para que um morfismo seja um monomorfismo, respectivamente um isomorfismo, é necessário e suficiente que seja radicial e não ramificado, respectivamente radicial, étale e sobrejetivo; daí a equivalência entre (ii) e (iii), respectivamente entre (ii') e (iii'). Resta provar que (i) equivale a (ii) e (i') a (ii').

Para isso, utiliza-se o critério diferencial de não ramificação (SGA I 3) e o isomorfismo
\[
\Omega^1_{X/S}\simeq \Omega^1_{X/X^{(p)}}
\]
dado pela sequência exata canônica
\[
\Fr_{X/S}^*(\Omega^1_{X^{(p)}/S})\longrightarrow \Omega^1_{X/S}
\longrightarrow \Omega^1_{X/X^{(p)}}\longrightarrow 0
\]
(SGA II 4), na qual a primeira flecha é visivelmente nula, pois $\Omega^1_{X^{(p)}/S}$ é gerada pelos diferenciais das seções de $\OO_{X^{(p)}}$, e o diferencial $d_{X/S}$ se anula sobre $(\OO_X)^p$ e sobre $\OO_S$. Assim, (i) equivale a (ii).
